\documentclass[../../main.tex]{subfiles}% 注意这里的文件路径不能用 ./main.tex ,否则用latexmk编译子文件会报错
\graphicspath{{\subfix{./image/}}{\subfix{../../image/}}} % 指定图片引用路径,后续可以直接使用图片文件名
% 如果图片存放在上级目录,则使用 ../../image/ 作为备用路径

% 例如:
% \begin{figure}[H]
% \centering
% \includegraphics[width=0.3\textwidth]{图.png}
% \caption{}
% \label{figure:图}
% \end{figure}
% 注意:上述\label{}一定要放在\caption{}之后,否则引用图片序号会只会显示??.

\begin{document}

\section{预备知识}\label{section:ch01-预备知识}

我们将使用集合论的通常记号。如通常，$\mathbb{N}$、$\mathbb{Z}$、$\mathbb{Q}$、$\mathbb{R}$ 和 $\mathbb{R}^N$（$N\in\mathbb{N}$；$\mathbb{R}^1\equiv\mathbb{R}$）分别表示自然数集、整数集、有理数集、实数集以及 $N$ 元实数组的集合。区间的熟知记号 $(a,b)$、$[a,b]$、$(a,b]$、$[a,b)$ 也将用于扩充实直线 $\overline{\mathbb{R}}\triangleq \mathbb{R}\cup\{\pm\infty\}$ 的子区间。偶尔我们使用符号 $[-\infty,+\infty]$ 表示 $\overline{\mathbb{R}}$；我们也用 $\infty$ 代替 $+\infty$。

\begin{definition}
设 $X$ 为一集合，定义其\textbf{幂集}为$\mathcal{P}(X)\triangleq \{E\mid E\subseteq X\}$。

若对所有 $i,j\in I$、$i\neq j$ 有 $F_i\cap F_j=\varnothing $。一个有限不交集族 $\{F_1,\cdots,F_n\}\subseteq\mathcal{P}(X)$ 若满足 $X=\bigcup_{i=1}^n F_i$，则称为 $X$ 的一个\textbf{划分}。
\end{definition}


\end{document}